% دستوراتی برای به حالت عادی درآمدن اندازه فونت‌ها و فاصله بین خطوط
\normalsize
\small
% مراجع خود را در این قسمت وارد کنید
%\begin{thebibliography}{99}
% چنانچه از نسخه‌های قدیمی زی‌پرشین استفاده می‌کنید، هنگام استفاده از دستورهای 

\begin{thebibliography}{99}

\bibitem{0}
دیوید کینکید، وارد چنی، آنالیز عددی، جلد دوم، مترجمان: دکتر فائزه توتونیان، منصوره صائمی، به‌نشر ۱۳۸۱.
\bibitem{1} 
ایند آر پینچ، کنترل بهینه و حساب تغییرات، مترجم: دکتر محمد هادی فراهی، موسسه چاپ و انتشارات دانشگاه فردوسی مشهد، 1385.

\bibitem{2}
دیوید برگز، الکساندر گراهام، مقدمه‌ای بر نظریه کنترل و کنترل بهینه، مترجمان دکتر علی وحیدیان کامیاد، دکتر ابوالقاسم بزرگ‌نیا، موسسه چاپ و انتشارات دانشگاه فردوسی مشهد، 1372.
%\setLTRbibitems
\begin{LTR}
\resetlatinfont
%  با خطا مواجه می‌شوید. به همین خاطر، باید از دستورهای زیر استفاده کنید.
% چنانچه مرجع فارسی هم دارید باید دستور Persian را فعال کرده و مراجع فارسی خود را بعد از این دستور وارد کنید
%\Persian
%\Latin

\bibitem{abramsky1}
S. Abramsky, {\em Domain theory in logical form}, Ann. Pure Applied Logic 51 (1991) 1–77.

\bibitem{abramsky2}
S. Abramsky, A. Jung, {\em Domain theory}, in: S. Abramsky, D.M. Gabbay, T.S.E. Maibaum (Eds.), Handbook of
Logic in Computer Science, Vol. 3, Clarendon Press, Oxford, 1994, pp. 1–68.

\bibitem{aliprantis}
C.D. Aliprantis and O. Burkinshaw,   {\em Principles of Real Analysis}.  Academic Press.

\bibitem{alvarez1}
M. Alvarez-Manilla, {\em Measure theoretic results for continuous valuations on partially ordered spaces}, Ph.D.
thesis, Imperial College, University of London, 2001.

\bibitem{alvarez2}
M. Alvarez-Manilla, A. Edalat, N. Saheb-Djahromi, {\em An extension result for continuous valuations}, J. London
Math. Soc. 61 (2000) 629–640.

\bibitem{mainarticle}
M. Alvarez-Manilla, A. Jung, K. Keimel, {\em The probabilistic powerdomain for stably compact
spaces}, Theoretical Computer Science 328 (2004) 221 – 244.

\bibitem{birkhoff}
G. Birkhoff, {\em Lattice Theory}, 3rd Edition, AMS Colloq. Publication, Vol. 25, American Mathematical Society,
Providence, 1967.

%\bibitem{choq1}
%G. Choquet, {\em Theory of capacities}, Ann. Inst. Fourier, Grenoble 5 (1953) 131–296.

\bibitem{choq3}
G. Choquet, {\em Lectures on Analysis}, Vol. 1, W. A. Benjamin Inc., London, 1969.

%\bibitem{choq2}
%G. Choquet, {\em Lectures on Analysis}, Vol. 2, W. A. Benjamin Inc., London, 1969.

\bibitem{desh}
J. Desharnais, V. Gupta, R. Jagadeesan, P. Panangaden, {\em Metrics for labeled Markov systems}, in: J.C.M.
Baeten, S. Mauw (Eds.), Proc. 10th Internat. Conf. on Concurrency Theory, Lecture Notes in Computer
Science, Vol. 1664, Springer, Berlin, 1999, pp. 258–273.

\bibitem{edward1}
D.A. Edwards, {\em On the existence of probability measures with given marginals}, Ann. Inst. Fourier, Grenoble,
28 (1978) 53–78.

\bibitem{folland}
G.B. Folland, {\em Real Analysis: Modern Techniques and Their Applications}, 2nd Edition, Wiley, 1999.

\bibitem{gierz1}
G. Gierz, K.H. Hofmann, K. Keimel, J.D. Lawson, M. Mislove, D.S. Scott, {\em A Compendium of Continuous
Lattices}, Springer, Berlin, 1980.

\bibitem{gierz2}
G. Gierz, K.H. Hofmann, K. Keimel, J.D. Lawson, M. Mislove, D.S. Scott, {\em Continuous Lattices and
Domains}, Encyclopedia of Mathematics and its Applications, Vol. 93, Cambridge University Press,
Cambridge, 2003.

%\bibitem{hofmann}
%K.H. Hofmann, {\em Stably continuous frames and their topological manifestations}, H.L. Bentley, H. Herrlich,
%M. Rajagopalan, H.Wolff (Eds.), Categorical Topology, Proc. 1983 Conf. Toledo, Sigma Series in Pure and
%AppliedMathematics, Vol. 5, Heldermann, Berlin, 1984, pp. 282–307.

\bibitem{horn}
A. Horn, A. Tarski, {\em Measures on Boolean algebras}, Trans. Amer. Math. Soc. 64 (1948) 467–497.

\bibitem{jones1}
C. Jones, {\em Probabilistic non-determinism}, Ph.D. thesis, University of Edinburgh, Edinburgh, 1990. Also
publishedas Technical Report No. CST-63-90.

\bibitem{jones2}
C. Jones, G. Plotkin, {\em A probabilistic powerdomain of evaluations}, in: Proc. 4th Annu. Symp. on Logic in
Computer Science, IEEE Computer Society Press, 1989, pp. 186–195.

\bibitem{jung1}
A. Jung, {\em Cartesian Closed Categories of Domains}, CWI Tracts, Vol. 66, Centrum voor Wiskunde en
Informatica, Amsterdam, 1989. 107 pp.

\bibitem{jung2}
A. Jung, {\em The classification of continuous domains}, in: Proc. 5th Annu. IEEE Symp. Logic in Computer
Science, IEEE Computer Society Press, Silverspring, MD, 1990, pp. 35–40.

%\bibitem{jung3}
%A. Jung, {\em Stably compact spaces and the probabilistic powerspace construction}, in: J. Desharnais, P.
%Panangaden (Eds.), Proc. Workshop on Domain-theoretic Methods in Probabilistic Processes, Electronic
%Notes in Theoretical Computer Science, Vol. 87, Elsevier Science Publishers B.V., Amsterdam, 2004.

\bibitem{jung4}
A. Jung, R. Tix, {\em The troublesome probabilistic powerdomain}, in: A. Edalat, A. Jung, K. Keimel, M.
Kwiatkowska (Eds.), Proc. 3rd Workshop on Computation and Approximation, Electronic Notes in
Theoretical Computer Science, Vol. 13, Elsevier Science Publishers B.V., Amsterdam, 1998, p. 23 pages.

%\bibitem{kegelmann}
%M. Kegelmann, {\em Continuous domains in logical form}, Ph.D. thesis, School of Computer Science, The
%University of Birmingham, 1999.
%
%\bibitem{keimel}
%K. Keimel, {\em The probabilistic powerdomain for the upwards topology of compact ordered spaces},
%J. Desharnais, P. Panangaden (Eds.), Proc. Workshop on Domain-theoretic Methods in Probabilistic
%Processes, Electronic Notes in Theoretical Computer Science, Vol. 87, Elsevier Science Publishers B.V.,
%Amsterdam, 2004.

\bibitem{kirch}
O. Kirch, {\em Bereiche und Bewertungen}. Master’s thesis, Technische Hochschule Darmstadt, June 1993, 77pp.

%\bibitem{Konig}
%H. Konig,  {\em Measure and Integration}, Springer–Verlag, 1997, xxi+260 pp.

\bibitem{kozen}
D. Kozen, {\em Semantics of probabilistic programs}, J. Comput. System Sci. 22 (1981) 328–350.

\bibitem{lawson1}
J.D. Lawson, {\em Valuations on continuous lattices}, in: R.E. Hoffmann (Ed.), Continuous Lattices and Related
Topics, Mathematik Arbeitspapiere, Vol. 27, Universität Bremen, 1982, pp. 204–225.

\bibitem{lawson2}
J.D. Lawson, {\em The versatile continuous order}, in: M. Main, A. Melton, M. Mislove, D. Schmidt (Eds.),
Mathematical Foundations of Programming Language Semantics, Lecture Notes in Computer Science,
Vol. 298, Springer, Berlin, 1988, pp. 134–160.

\bibitem{lawson3}
J.D. Lawson, {\em Domains, Integration, and Positive Analysis}. Mathematical Structures in
Computer Science, 14 , 2004, 815-832.

\bibitem{moshier}
M.A. Moshier, A. Jung, {\em A logic for probabilities in semantics}, in: J. Bradfield (Ed.), Computer Science Logic,
Lecture Notes in Computer Science, Vol. 2471, Springer, Berlin, 2002, pp. 216–231.

%\bibitem{mun}
%J.R. Munkres, {\em Topology: A First Course}, Prentice-Hall, 1975.

\bibitem{rudin} 
W. Rudin,  {\em Real and Complex Analysis}.  Mc Graw-Hill Book Comp. 1966.

\bibitem{functional}
W. Rudin,  {\em Functional Analysis},  2nd Edition, Mc Graw-Hill Book Comp.  1991.

\bibitem{saheb}
N. Saheb-Djahromi, {\em CPO’s of measures for nondeterminism}, Theoret. Comput. Sci. 12 (1980) 19–37.

%\bibitem{scott1}
%D.S. Scott, {\em Outline of a mathematical theory of computation}, in: 4th Ann. Princeton Conf. Information
%Sciences and Systems, 1970, pp. 169–176.
%
%\bibitem{scott2}
%D.S. Scott, {\em Domains for denotational semantics}, in: M. Nielson, E.M. Schmidt (Eds.), International
%Colloquium on Automata, Languages and Programs, Lecture Notes in Computer Science,Vol. 140, Springer,
%Berlin, 1982, pp. 577–613.

\bibitem{smyth1}
M.B. Smyth, {\em Powerdomains and predicate transformers}: a topological view, in: J. Diaz (Ed.), Automata,
Languages andProgramming, Lecture Notes in Computer Science, Vol. 154, Springer, Berlin, 1983,
pp. 662–675.

\bibitem{smyth2}
M.B. Smyth, \textit{Topology}, S. Abramsky, D.M. Gabbay, T.S.E. Maibaum (Eds.), {\em Handbook of Logic in
Computer Science}, Vol. 1, Clarendon Press, Oxford, 1992, pp. 641–761.

%\bibitem{tix}
%R. Tix, {\em Stetige Bewertungen auf topologischen Räumen}, Master’s thesis, Technische Hochschule Darmstadt,
%June 1995, 51pp.

\bibitem{topsze}
F. Topsze, {\em Topology and Measure}, Lecture Notes in Mathematics, Vol. 133, Springer, Berlin, 1970.

\bibitem{vickers}
S.J. Vickers, {\em Topology Via Logic}, Cambridge Tracts in Theoretical Computer Science, Vol. 5, Cambridge
University Press, Cambridge, 1989.

\end{LTR}
\end{thebibliography}